\section{The a priori bounds}
\label{sec:bnd}

Every truncation in this document is controlled by a proved bound, and every bound is
checked numerically in \code{BigFloat} against the quantity it bounds. This section states
the bounds, proves them, says exactly what each one bounds and relative to what scale, and
gives the measured ratio of bound to truth.

\begin{quote}
\textbf{Which bound the shipped code uses.} The $l$-truncation bound of this section,
Theorem~\ref{thm:bnd:rem}, is \emph{not} the one \code{farfield.jl} evaluates.
Section~\ref{sec:bnd2} replaces it by two sharper theorems, and it is those --- Theorem~A
(Theorem~\ref{thm:bnd2:reg}, whole box) and Theorem~B (Theorem~\ref{thm:bnd2:sub},
sub-box) --- that select
$L$ in the library, at the tolerance $\mathrm{tol} = 10^{-14}$
(Section~\ref{sec:join:rule}). Theorem~\ref{thm:bnd:rem} is kept here in full because it is
the bound every measurement of the first round was made against, because
Theorem~\ref{thm:bnd2:sng} is its sharpened form, and because the two bound
\emph{different} truncations: Theorem~\ref{thm:bnd:rem} bounds the sum with the derivatives
on the singular side, and the library truncates the one with the derivatives on the regular
side (Section~\ref{sec:bnd2:which}). Everything else in this section --- the elementary
bounds on $j_l$ and $h_l$, the $n$-truncation, the $k$-series bound and the two radii ---
is used unchanged.
\end{quote}

\noindent
The scripts are
\code{chkElem.jl} (the elementary bounds (a), (b), the scalar addition-theorem
remainder, the radial moments $W_l$, and the $n$-tail (d)), \code{chk17.jl} (the two
structural constants of (c)), \code{chkRem.jl} (the $l$-truncation bound end to end),
\code{chkKser.jl} (the $k$-series bound (e)) and \code{chkRay.jl} (the convergence radius
of remark (f)); their raw output is \code{out\_elem.txt}, \code{out\_c17.txt},
\code{out\_rem.txt}, \code{out\_kser.txt}, \code{out\_ray.txt}.

\subsection{What is bounded, and the scale it is measured against}
\label{sec:bnd:scale}

The bounds below are \emph{absolute} bounds on the truncation error of $T_{ab}$. A
relative statement needs a scale, and using $\max_{ab}|T_{ab}|$ would be circular, since
the tensor is what is being computed. The scale used throughout is the a priori magnitude
\begin{equation}
\mathrm{est}(\mathbf{R}) \;=\;
\frac{V_{\mathrm{t}}\,|k|^2\,e^{-\operatorname{Im}(k)|\mathbf{R}|}}
     {4\pi|f|^2|\mathbf{R}|}
\left(1 + \frac{3}{|k\mathbf{R}|} + \frac{3}{|k\mathbf{R}|^2}\right),
\label{eq:bnd:est}
\end{equation}
the pointwise dyadic scale $|(\partial_a\partial_b+\delta_{ab}k^2)g|(\mathbf{R})$
majorised term by term, multiplied by the total weight
$(1/V_{\mathrm{t}})\int_D w = V_{\mathrm{t}}$. It is computable before anything else is
and it is never optimistic: measured against a $220$-bit reference tensor,
$\mathrm{est}/\max_{ab}|T_{ab}| = 2.36$, $2.26$, $2.59$, $3.77$, $1.94$ and $2.29$ at
$(4,0,0)$ and $(3,1,0)$ of the $\lambda/32$ cube (both frequencies), $(8,8,8)$ of the
$\lambda/32$ cube, $(4,0,0)$ of the $\lambda/4$ cube and $(4,0,0)$ of the slender cell
(\code{chkRem.jl}), and $2.0$ to $5.1$ over the wider set of offsets and shapes on which it
was measured elsewhere. A tolerance of $10^{-13}$ against $\mathrm{est}$ is therefore
between $2$ and $5$ times stricter than $10^{-13}$ against the largest entry.

\subsection{The two elementary bounds}
\label{sec:bnd:elem}

\begin{theorem}[regular factor]
\label{thm:bnd:j}
For every $l \ge 0$ and every $z \in \mathbb{C}$,
\begin{equation}
|j_l(z)| \;\le\; \frac{|z|^{l}}{(2l+1)!!}\;e^{|z|^{2}/(4l+6)} .
\label{eq:bnd:j}
\end{equation}
\end{theorem}

\begin{proof}
$j_l(z) = \dfrac{z^l}{(2l+1)!!}\displaystyle\sum_{n\ge0}
\frac{(-z^2/2)^n}{n!\prod_{q=1}^{n}(2l+2q+1)}$. Since $2l+2q+1 \ge 2l+3$ for every
$q \ge 1$, $\prod_{q=1}^{n}(2l+2q+1) \ge (2l+3)^n$, so by the triangle inequality
\[
|j_l(z)| \le \frac{|z|^l}{(2l+1)!!}\sum_{n\ge0}\frac{(|z|^2/2)^n}{n!\,(2l+3)^n}
= \frac{|z|^l}{(2l+1)!!}\,e^{|z|^2/(2(2l+3))},
\]
and $2(2l+3) = 4l+6$.
\end{proof}

\begin{theorem}[singular factor]
\label{thm:bnd:h}
For every $l \ge 0$ and every $z \neq 0$,
\begin{equation}
\bigl|h^{(1)}_l(z)\bigr| \;\le\; \mathrm{hb}(l,z) \;:=\;
\frac{e^{-\operatorname{Im}z}}{|z|}\sum_{s=0}^{l}
\frac{(l+s)!}{s!\,(l-s)!\,(2|z|)^{s}} .
\label{eq:bnd:h}
\end{equation}
Moreover $\mathrm{hb}(l,z)$ is increasing in $l$ at fixed $z$.
\end{theorem}

\begin{proof}
The classical finite representation is
$h^{(1)}_l(z) = (-\ii)^{l+1}\dfrac{e^{\ii z}}{z}\sum_{s=0}^{l}\ii^{s}
\dfrac{(l+s)!}{s!(l-s)!(2z)^{s}}$, and $|e^{\ii z}| = e^{-\operatorname{Im}z}$; apply the
triangle inequality. Monotonicity in $l$ holds term by term because
$(l+s)!/(s!(l-s)!)$ increases with $l$ at fixed $s$ and the sum gains a term.
\end{proof}

\paragraph{Measured (script \code{chkElem.jl}, $400$-bit).} Over $64$ rows of
\eqref{eq:bnd:j} --- $l \in \{0,1,2,5,10,20,40,80\}$ and $z \in \{0.3, 1, 3, 10, 30,
3+0.3\ii, 10+3\ii, 30+30\ii\}$ --- the bound is never violated and
$\mathrm{bound}/|j_l|$ ranges from $1.0011$ to $1.5\times10^{119}$. The tightness is
entirely a function of $|z|^2/(2l+3)$, which is the argument regime that matters: at
$|z| = 0.3$ the ratio is $1.031$, $1.0039$, $1.0011$, $1.00055$ at $l = 0, 10, 40, 80$; at
$|z| = 1$ it is $1.404$, $1.044$, $1.012$, $1.0062$; at $|z| = 3$ it is $95.3$ (at
$l = 0$, where $j_0$ sits near its zero at $z = \pi$), $1.48$, $1.11$, $1.057$. Family
(b) evaluates $j_l$ only at $|z| = |k|\,|\boldsymbol{\delta}| \le |k| r_{\mathrm{d}}$,
which is $0.34$ for the $\lambda/32$ cube, $1.36$ at $\lambda/8$, $2.72$ at $\lambda/4$
and $0.278$ for the slender cell, so \eqref{eq:bnd:j} is tight to within a factor
$1.06$--$1.5$ everywhere it is used. The enormous ratios in the table occur only for
$|z|^2 \gg l$, where the bound's $e^{|z|^2/(4l+6)}$ is the wrong majorant --- a regime
this expansion never enters.

Over $72$ rows of \eqref{eq:bnd:h} the bound is never violated, with equality at $l = 0$
(both sides are $e^{-\operatorname{Im}z}/|z|$), and $\mathrm{bound}/|h_l|$ runs from
$1.0$ to $1.8\times10^{13}$. Again the loose rows are a single regime, $|z| \lesssim l$
(the terms $(l+s)!/(s!(l-s)!(2|z|)^s)$ grow before they decay, so the triangle inequality
gives away the oscillation): the ratio is $1.35$ at $|z| = 0.3$, $l = 80$, but
$1.2\times10^{4}$ at $|z| = 10$, $l = 80$, and $3.0\times10^{3}$ at $|z| = 100$,
$l = 40$. That loss is one of the two main reasons the composite bound of
Theorem~\ref{thm:bnd:rem} over-selects $L$.

\subsection{\texorpdfstring{The $l$-truncation bound}{The l-truncation bound}}
\label{sec:bnd:rem}

\begin{theorem}[$l$-truncation of \eqref{eq:exp:eval}]
\label{thm:bnd:rem}
Let the pair be separated and $r_{\mathrm{d}} = |\mathbf{s}|$. Write
$T^{(L)}_{ab}$ for the sum \eqref{eq:exp:eval} truncated at $l \le L$. Then for every
$a,b$
\begin{equation}
\bigl|T_{ab}(\mathbf{R}) - T^{(L)}_{ab}(\mathbf{R})\bigr|
\;\le\;
\frac{17\,|k|^{3}}{4\pi|f|^{2}V_{\mathrm{t}}}
\sum_{\substack{l > L\\ l \text{ even}}}
(2l+1)\sqrt{2l+5}\;
\frac{|k|^{l}\,e^{(|k|r_{\mathrm{d}})^{2}/(4l+6)}}{(2l+1)!!}\;
W_l\;\mathrm{hb}\bigl(l+2,\,k|\mathbf{R}|\bigr),
\label{eq:bnd:rem}
\end{equation}
where $W_l = \int_D w(\boldsymbol{\delta})|\boldsymbol{\delta}|^{l}\,
\dd\boldsymbol{\delta}$ is, for even $l$, the exact positive sum
\begin{equation}
W_{2p} = \sum_{i+j+q=p}\frac{p!}{i!\,j!\,q!}\,
\mu_{2i}(s_1)\,\mu_{2j}(s_2)\,\mu_{2q}(s_3),
\label{eq:bnd:W}
\end{equation}
and $\mathrm{hb}$ is the majorant \eqref{eq:bnd:h}.
\end{theorem}

\begin{proof}
Start from \eqref{eq:exp:sing}, the organization with the derivatives on the singular
side. Three estimates are needed.

\emph{(i) The derivative factor, and the constant $17$.} Each Cartesian first derivative
of a singular solid harmonic $h_l(k|\mathbf{R}|)Y_{lm}(\hat{\mathbf{R}})$ is $k$ times a
combination of at most four terms $h_{l\pm1}Y_{l\pm1,\,m\text{ or }m\pm1}$ with
coefficients of modulus at most one; these are the standard differentiation relations, in
which the coefficients are of the form
$a_{lm} = \sqrt{(l-m)(l+m)/((2l-1)(2l+1))} \le 1$ and its companions. Two derivatives
therefore give at most $16$ terms with $l' \in \{l-2,l,l+2\}$, and the term
$\delta_{ab}k^2 h_l Y_{lm}$ adds one more: at most $17$ terms, each of modulus at most
$|k|^2 \max_{l'}\mathrm{hb}(l',|k\mathbf{R}|)\max_{m'}|Y_{l'm'}|$. By
Theorem~\ref{thm:bnd:h}, $\mathrm{hb}$ is increasing in $l'$, so the maximum is at
$l' = l+2$; and $|Y_{l'm'}| \le \sqrt{(2l'+1)/(4\pi)} = \sqrt{(2l+5)/(4\pi)}$. Hence
\begin{equation}
\Bigl|(\partial_a\partial_b+\delta_{ab}k^2)
\bigl(h_l(k|\mathbf{R}|)Y_{lm}(\hat{\mathbf{R}})\bigr)\Bigr|
\;\le\; 17\,|k|^{2}\,\mathrm{hb}(l+2,k|\mathbf{R}|)\,\sqrt{\frac{2l+5}{4\pi}} .
\label{eq:bnd:17}
\end{equation}

\emph{(ii) The sum over $m$.} By Cauchy--Schwarz and the addition theorem at coincident
arguments, $\sum_{m}Y_{lm}(\mathbf{u})^2 = (2l+1)P_l(1)/(4\pi) = (2l+1)/(4\pi)$, so
\begin{equation}
\sum_{m=-l}^{l}\bigl|Y_{lm}(\mathbf{u})\bigr|
\;\le\; \sqrt{(2l+1)\sum_m Y_{lm}(\mathbf{u})^2}
\;=\; \frac{2l+1}{\sqrt{4\pi}} .
\label{eq:bnd:Ysum}
\end{equation}

\emph{(iii) The regular factor and the radial moment.} By Theorem~\ref{thm:bnd:j}
pointwise in $\boldsymbol{\delta}$, and since
$|\boldsymbol{\delta}| \le r_{\mathrm{d}}$ on $D$,
\[
\bigl|j_l(k|\boldsymbol{\delta}|)\bigr| \le
\frac{|k|^{l}|\boldsymbol{\delta}|^{l}}{(2l+1)!!}\,
e^{(|k|r_{\mathrm{d}})^{2}/(4l+6)} ,
\]
so, applying \eqref{eq:bnd:Ysum} to $\hat{\boldsymbol{\delta}}$ under the integral and
using $w \ge 0$,
\[
\sum_m \bigl|\mathcal{I}_{lm}\bigr|
\le \frac{2l+1}{\sqrt{4\pi}}\,
\frac{|k|^{l}e^{(|k|r_{\mathrm{d}})^{2}/(4l+6)}}{(2l+1)!!}
\int_D w\,|\boldsymbol{\delta}|^{l}\,\dd\boldsymbol{\delta}
= \frac{2l+1}{\sqrt{4\pi}}\,
\frac{|k|^{l}e^{(|k|r_{\mathrm{d}})^{2}/(4l+6)}}{(2l+1)!!}\,W_l .
\]
Multiplying (i) by this, collecting the prefactor $|k|/(|f|^2V_{\mathrm{t}})$ of
\eqref{eq:exp:sing} and the two factors $1/\sqrt{4\pi}$ into $1/(4\pi)$ gives
\eqref{eq:bnd:rem}; the sum runs over even $l$ only because
Section~\ref{sec:exp:parity} shows $\mathcal{I}_{lm} = 0$ for odd $l$. Finally
\eqref{eq:bnd:W} is the multinomial expansion of $|\boldsymbol{\delta}|^{2p} =
(\delta_1^2+\delta_2^2+\delta_3^2)^p$ integrated term by term against
\eqref{eq:vol:mu}; every term is positive, so $W_l$ is computed without cancellation.
\end{proof}

\begin{remark}[why $W_l$ and not $V_{\mathrm{t}}^2 r_{\mathrm{d}}^{\,l}$]
\label{rem:bnd:W}
Bounding $|\boldsymbol{\delta}|^l$ by $r_{\mathrm{d}}^{\,l}$ under the integral is legal
and gives $W_l \le V_{\mathrm{t}}^{2}r_{\mathrm{d}}^{\,l}$, but it throws away exactly the
feature that makes the expansion converge faster than $\rho^l$: the corner of $D$, where
$|\boldsymbol{\delta}| = r_{\mathrm{d}}$, is where the triangle weight vanishes to third
order. With the crude majorant the measured bound-to-truth ratio is $10^{5}$ to
$10^{12}$; with the exact $W_l$ it is $10^{3}$ to $10^{7}$ (below). The practical
consequence is that the effective convergence ratio for a cube is $0.68$ to $0.71$ times
$\rho$ rather than $\rho$, so the empirical rule for the cube is
$L(10^{-13}) \approx 13/\log_{10}(1/(0.7\rho))$.
\end{remark}

\paragraph{The two constants, measured (script \code{chk17.jl}).} \eqref{eq:bnd:Ysum} is
never violated over $l \le 60$ and four directions; it is an equality at $l = 0$ and the
ratio settles to a direction-dependent constant, $0.487$ along an axis, $0.717$ on the
body diagonal, $0.654$ on a generic direction at $l = 60$, so it costs at most a factor
$2$ and not a factor growing with $l$. \eqref{eq:bnd:17} is far looser: over $72$ rows
($\mathbf{n} \in \{(2,0,0),(4,0,0),(3,1,0),(16,16,16)\}$ at $\lambda/32$ and
$\{(4,0,0),(32,0,0)\}$ at $\lambda/4$, $f = 1$ and $1+0.1\ii$, $l \le 24$, all nine
$(a,b)$ and all $m$, derivatives taken by $320$-bit central differences at step
$10^{-25}$) the worst ratio of the true derivative to the right-hand side of
\eqref{eq:bnd:17} is $3.14\times10^{-2}$ and the best is $2.94\times10^{-5}$. \emph{The
constant $17$ is loose by a factor of $32$ to $34000$}, which is where most of the
over-selection of Theorem~\ref{thm:bnd:rem} comes from: it counts $17$ terms that could
all be maximal at once, whereas the differentiation relations distribute the weight
across $l' = l-2, l, l+2$ and across $m$.

\paragraph{The bound end to end (script \code{chkRem.jl}, $320$-bit).} The truth is
$\max_{ab}|T_{ab} - T^{(L)}_{ab}|$ with $T_{ab}$ from an independent \code{BigFloat}
tensor Gauss--Legendre evaluation of \eqref{eq:vol:main} on the eight octants of $D$ at
order $26$ (self-consistency against order $20$: $10^{-37}$ to $10^{-67}$ depending on
the offset), and $T^{(L)}$ assembled from \eqref{eq:exp:sing} with $\mathcal{I}_{lm}$ by
an order-$14$ octant rule and the derivatives by central differences. Over $79$ rows the
bound is never violated:

\begin{center}\footnotesize
\begin{tabular}{@{}llrrrrrr@{}}
\toprule
shape & $f$ & $\mathbf{n}$ & $\rho$ & $|k\mathbf{R}|$ & $L$
& $\max_{ab}|\mathrm{Rem}_L|$ & bound/truth\\
\midrule
$(1/32)^3$ & $1$        & $(4,0,0)$    & 0.433 & 0.785 & 8  & 1.62e-9  & 3.17e4\\
$(1/32)^3$ & $1$        & $(4,0,0)$    & 0.433 & 0.785 & 16 & 2.53e-13 & 5.50e4\\
$(1/32)^3$ & $1$        & $(4,0,0)$    & 0.433 & 0.785 & 24 & 2.42e-17 & 2.21e5\\
$(1/32)^3$ & $1+0.1\ii$ & $(4,0,0)$    & 0.433 & 0.789 & 16 & 2.50e-13 & 5.10e4\\
$(1/32)^3$ & $1$        & $(3,1,0)$    & 0.548 & 0.621 & 16 & 1.16e-11 & 1.57e5\\
$(1/32)^3$ & $1$        & $(3,1,0)$    & 0.548 & 0.621 & 24 & 7.09e-14 & 6.54e4\\
$(1/32)^3$ & $1$        & $(8,8,8)$    & 0.125 & 2.721 & 8  & 1.29e-16 & 2.04e5\\
$(1/32)^3$ & $1$        & $(8,8,8)$    & 0.125 & 2.721 & 16 & 4.23e-25 & 8.51e5\\
$(1/4)^3$  & $1$        & $(4,0,0)$    & 0.433 & 6.283 & 8  & 4.19e-9  & 1.57e6\\
$(1/4)^3$  & $1$        & $(4,0,0)$    & 0.433 & 6.283 & 16 & 4.17e-13 & 5.50e6\\
$(1/4)^3$  & $1$        & $(4,0,0)$    & 0.433 & 6.283 & 24 & 3.31e-17 & 2.99e7\\
sl         & $1$        & $(4,0,0)$    & 0.354 & 0.785 & 8  & 4.44e-11 & 1.91e4\\
sl         & $1$        & $(4,0,0)$    & 0.354 & 0.785 & 16 & 1.62e-15 & 5.93e4\\
sl         & $1$        & $(4,0,0)$    & 0.354 & 0.785 & 24 & 1.49e-19 & 9.47e4\\
\bottomrule
\end{tabular}
\end{center}

\noindent
The full range over the $79$ rows is $\mathrm{bound}/\mathrm{truth} \in
[1.3\times10^{3},\,3.0\times10^{7}]$, the largest values at $\lambda/4$ where
$|k\mathbf{R}| = 6.3$ and the $\mathrm{hb}$ majorant of Theorem~\ref{thm:bnd:h} is at its
loosest. In terms of the quantity that matters, the selected $L$, an over-estimate of
$10^{3}$ to $10^{7}$ in the bound costs a factor $\log_{10}(10^{3\dots7})/13 \approx
0.23$ to $0.54$ in $L(10^{-13})$, i.e.\ the bound asks for $1.2$ to $1.5$ times the $L$
that is actually needed. \emph{The bound is a usable selector, not merely a certificate},
and Section~\ref{sec:join} uses it as one.

\subsection{The sub-box version}
\label{sec:bnd:sub}

Theorem~\ref{thm:bnd:rem} applies verbatim to each sub-box of \eqref{eq:exp:split} with
three changes: the sum runs over \emph{all} $l$, because subdivision destroys the parity
of Section~\ref{sec:exp:parity}; the expansion centre is $\mathbf{V}_j =
\mathbf{R}+\mathbf{c}_j$; and the radial moment is the sub-box's own,
\begin{equation}
\bigl|T^{(j)}_{ab}-T^{(j),L_j}_{ab}\bigr| \le
\frac{17|k|^{3}}{4\pi|f|^{2}V_{\mathrm{t}}}\sum_{l>L_j}(2l+1)\sqrt{2l+5}
\frac{|k|^{l}e^{(|k|r_j)^{2}/(4l+6)}}{(2l+1)!!}\,W^{(j)}_l\,
\mathrm{hb}\bigl(l+2, k|\mathbf{V}_j|\bigr),
\label{eq:bnd:sub}
\end{equation}
$r_j = |\mathbf{h}_j|$, with
$W^{(j)}_{2p} = \sum_{i+j'+q=p}\frac{p!}{i!j'!q!}
\nu_{2i}(\alpha_1,\cdot,h_1)\nu_{2j'}(\alpha_2,\cdot,h_2)\nu_{2q}(\alpha_3,\cdot,h_3)$
--- only even one-dimensional moments enter, so $\beta$ drops out and every term is
positive --- and $W^{(j)}_l \le \sqrt{W^{(j)}_{l-1}W^{(j)}_{l+1}}$ for odd $l$ by
Cauchy--Schwarz with the non-negative weight. Each of the $n_{\mathrm{sub}}$ sub-boxes is
given $\mathrm{tol}\cdot\mathrm{est}(\mathbf{R})/n_{\mathrm{sub}}$ of the budget and
$L_j$ is the first $l$ at which its tail falls below it. The sub-box bound is
\emph{tighter in $L$} than the whole-box one, by $1.2$--$1.3$ against $1.5$--$2.0$,
because $W^{(j)}_l$ describes the affine weight of a sub-box far better than $W_l$
describes the corner suppression of the triangle weight.

\subsection{\texorpdfstring{The $n$-truncation of the $k$-series inside each $l$}{The n-truncation of the k-series inside each l}}
\label{sec:bnd:n}

\begin{theorem}[$n$-truncation]
\label{thm:bnd:n}
Fix $l$ and truncate the inner sum \eqref{eq:exp:cln} at $n \le N$. With
$y_l = (|k|r_{\mathrm{d}})^{2}/(2(2l+3))$, the omitted part of shell $l$ satisfies
\begin{equation}
\sum_{n>N}|c_{ln}(k)|\;W_{l+2n}
\;\le\;
\frac{|k|^{l}\,W_l}{(2l+1)!!}\;\frac{y_l^{\,N+1}}{(N+1)!}\;e^{y_l},
\label{eq:bnd:n}
\end{equation}
and the induced error in $T_{ab}$ is \eqref{eq:bnd:rem} with the sum restricted to
$l \le L$ and with $e^{(|k|r_{\mathrm{d}})^{2}/(4l+6)}$ replaced by
$y_l^{\,N+1}e^{y_l}/(N+1)!$.
\end{theorem}

\begin{proof}
$W_{l+2n} = \int_D w|\boldsymbol{\delta}|^{l+2n} \le r_{\mathrm{d}}^{2n}W_l$ because
$|\boldsymbol{\delta}| \le r_{\mathrm{d}}$ on $D$ and $w \ge 0$. From
\eqref{eq:exp:cln} and $(2l+2n+1)!! = (2l+1)!!\prod_{q=1}^{n}(2l+2q+1) \ge
(2l+1)!!\,(2l+3)^n$,
\[
|c_{ln}(k)|\,W_{l+2n} \le
\frac{|k|^{l}W_l}{(2l+1)!!}\cdot\frac{1}{n!}
\left(\frac{|k|^{2}r_{\mathrm{d}}^{2}}{2(2l+3)}\right)^{n}
= \frac{|k|^{l}W_l}{(2l+1)!!}\cdot\frac{y_l^{\,n}}{n!},
\]
and $\sum_{n>N}y^{n}/n! = (y^{N+1}/(N+1)!)\sum_{j\ge0}y^{j}(N+1)!/(N+1+j)! \le
(y^{N+1}/(N+1)!)\,e^{y}$ since $(N+1)!/(N+1+j)! \le 1/j!$. The second statement follows
by substituting this in place of the corresponding factor in step (iii) of the proof of
Theorem~\ref{thm:bnd:rem}, all other steps being unchanged.
\end{proof}

\paragraph{Measured (script \code{chkElem.jl}).} \eqref{eq:bnd:n} was checked in its
scalar form, $|j_l(z)-\sum_{n\le N}| \le \frac{|z|^{l}}{(2l+1)!!}
\frac{y^{N+1}}{(N+1)!}e^{y}$, at the four physical values
$|z| = |k|r_{\mathrm{d}} \in \{0.34, 1.36, 2.72, 0.278\}$ (the $\lambda/32$, $\lambda/8$
and $\lambda/4$ cubes and the slender cell), $l \in \{0,4,12,24\}$ and
$N \in \{2,4,6\}$: $48$ rows, no violation, bound/truth from $1.122$ to $3268$. The bound
is tight where it is used --- $1.12$ to $2.4$ for $l \ge 12$ at every $|z|$ --- and loose
only at $l = 0$ and large $N$, where $j_0$ is nearly its own leading term. It confirms the
term counts $n_{\max} = 6$ (at $\lambda/32$ and for the slender cell), $10$ (at
$\lambda/8$) and $12$ (at $\lambda/4$) for a relative truncation of $10^{-16}$, so the
$n$ series is a non-issue: the relative size of the $n$-th term of
$\widetilde{Q}_{lm}$, maximised over $l \le 24$, $m$ and entry type, is

\begin{center}\footnotesize
\begin{tabular}{@{}lrrrrrrr@{}}
\toprule
shape & $|k|r_{\mathrm{d}}$ & $n=0$ & $n=2$ & $n=4$ & $n=6$ & $n=8$ & $n=12$\\
\midrule
cube $\lambda/32$ & 0.340 & 2.0 & 2.5e-3 & 1.9e-9 & 4.6e-16 & 3.9e-23 & 3.1e-38\\
cube $\lambda/8$  & 1.360 & 2.1 & 4.2e-2 & 8.4e-6 & 5.1e-10 & 1.1e-14 & 5.8e-25\\
cube $\lambda/4$  & 2.721 & 2.4 & 2.0e-1 & 6.8e-4 & 6.4e-7  & 2.2e-10 & 2.9e-18\\
slender           & 0.278 & 2.0 & 1.8e-3 & 1.1e-9 & 1.7e-16 & 9.6e-24 & 3.3e-39\\
\bottomrule
\end{tabular}
\end{center}

\noindent
with its own cancellation $\sum_n|\text{term}|/|\text{sum}|$ equal to $3.07$, $3.29$ and
$4.14$ at $\lambda/32$, $\lambda/8$ and $\lambda/4$: only the $n=0$ entry exceeds one, so
the series alternates with a first term about twice the total.

\begin{remark}[$N$ is adaptive, and the library refuses rather than truncates]
\label{rem:bnd:adapt}
The term counts above are functions of $|k|r_{\mathrm{d}}$ alone --- a
\emph{shape and frequency} quantity, not a constant --- and the number they produce must be
built into the geometry table, because the table stores the first $N+1$ coefficients
$c_{ln}$ of \eqref{eq:exp:cln}. An earlier version of the library carried a compile-time cap
$n_{\max} = 12$ that the per-$l$ cut could only lower, so a shape needing more got $12$ and
no warning; the adversarial audit found the error growing as
$(|k|r_{\mathrm{d}})^{25.6}$ above $|k|r_{\mathrm{d}} \approx 3.2$, reaching
$9.6\times10^{-3}$ per entry at $|k|r_{\mathrm{d}} = 8.96$ with a reported bound
$7.2\times10^{8}$ times smaller (Section~\ref{sec:def}, defect D9). In the shipped library
the cut returned by Theorem~\ref{thm:bnd:n} is honoured: $N$ is taken as the largest
per-$l$ cut the bound asks for, the table is built at $\max(12, N)$, $N$ is re-checked at
every evaluation, and above $N_{\mathrm{cap}} = 40$ the call raises an error instead of
returning a wrong number. The measured relation is
\begin{center}\footnotesize
\begin{tabular}{@{}lrrrrrrrrrrr@{}}
\toprule
$|k|r_{\mathrm{d}}$ & 0.03 & 0.28 & 0.68 & 1.58 & 2.72 & 3.16 & 4.46 & 5.80 & 8.96 & 10.9 & 21.9\\
$N$ & 4 & 6 & 8 & 11 & 12 & 16 & 18 & 19 & 25 & 24 & 39\\
\bottomrule
\end{tabular}
\end{center}
\noindent
(70 rows, evaluated table-free at $\mathrm{tol} = 10^{-14}$; the non-monotone pairs are
different shapes at the same $|k|r_{\mathrm{d}}$, since $N$ also depends weakly on the
aspect ratio). A cell whose longest edge is at most $\lambda/4$ has
$r_{\mathrm{d}} \le \sqrt3\,s_{\max}$ and hence $|k|r_{\mathrm{d}} \le 2.72$, so
$N \le 12$: \emph{the old constant was exactly right for Gila's stated range and wrong for
anything coarser}, which is why the cap was invisible until the audit went outside it.
$N_{\mathrm{cap}} = 40$ covers $|k|r_{\mathrm{d}} \le 22$, a cubic cell of one wavelength at
$f = 2$. On the four shapes Gila uses, the largest per-shell cut the bound asks for is
(\code{tables/e\_ncut.txt}, $\mathrm{tol} = 10^{-14}$)
\begin{center}\small
\begin{tabular}{@{}lrrrr@{}}
\toprule
shape & $(1/32)^3$ & $(1/8)^3$ & $(1/4)^3$ & $(1/32,1/32,1/512)$\\
\midrule
$|k|r_{\mathrm{d}}$ & $0.340$ & $1.360$ & $2.721$ & $0.278$\\
$N$ needed & 6 & 9 & 12 & 6\\
$N$ stored & 12 & 12 & 12 & 12\\
\bottomrule
\end{tabular}
\end{center}
\noindent
at both $f = 1$ and $f = 1+0.1\ii$, so the $n$-series is certified everywhere the library
evaluates on these shapes, with the cut falling from $6$ at $l = 0$ to $4$ at $l = 14$ on
the finest cell and from $12$ to $7$ on the coarsest.
\end{remark}

\subsection{\texorpdfstring{The $k$-series bound, for the near band}{The k-series bound, for the near band}}
\label{sec:bnd:kser}

The near field uses the moment series of the companion document, and the join needs its
bound in the same currency.

\begin{theorem}[$k$-series truncation]
\label{thm:bnd:kser}
For a panel pair $(A,B)$ write $I_m = \int_A\int_B|\mathbf{x}-\mathbf{y}|^{m}\dd S'\dd S$,
$D_{\max} = \max_{\mathbf{x}\in A,\mathbf{y}\in B}|\mathbf{x}-\mathbf{y}|$ and
$x = |k|D_{\max}$, and let
$S = \int_A\int_B g = \frac{1}{4\pi f^{2}}\sum_{n\ge0}\frac{(\ii k)^{n}I_{n-1}}{n!}$.
Then for every $N \ge 0$
\begin{equation}
\Bigl|S - \frac{1}{4\pi f^{2}}\sum_{n\le N}\frac{(\ii k)^{n}I_{n-1}}{n!}\Bigr|
\;\le\;
\frac{I_{-1}}{4\pi|f|^{2}}\;\frac{x^{N+1}}{(N+1)!}\;B(x,N),
\qquad
B(x,N) = \min\Bigl(e^{x},\ \frac{1}{1-x/(N+2)}\Bigr),
\label{eq:bnd:kser}
\end{equation}
the second branch only when $x < N+2$. Truncating every face pair at the same $N$ and
assembling,
\begin{equation}
\frac{\max_{ab}|G_{ab}-G^{(N)}_{ab}|}{\max_{ab}|G_{ab}|}
\;\le\; \Lambda\;\frac{x_{\max}^{N+1}}{(N+1)!}\,B(x_{\max},N),
\qquad
\Lambda = \frac{\sum_{\mathrm{fp}}I_{-1}^{(\mathrm{fp})}}
{4\pi|f|^{2}V_{\mathrm{t}}\max_{ab}|G_{ab}|} .
\label{eq:bnd:kasm}
\end{equation}
\end{theorem}

\begin{proof}
Pointwise on $A\times B$, $r^{\,n-1} = r^{-1}r^{n} \le r^{-1}D_{\max}^{n}$ and
$r^{-1} > 0$, so $|I_{n-1}| \le D_{\max}^{n}I_{-1}$ and $|t_n| \le x^{n}I_{-1}/n!$ with
$t_n = (\ii k)^{n}I_{n-1}/n!$. Then
$\sum_{n>N}x^{n}/n! = \frac{x^{N+1}}{(N+1)!}\sum_{j\ge0}x^{j}\frac{(N+1)!}{(N+1+j)!}$,
and $(N+1)!/(N+1+j)! \le 1/j!$ gives the $e^{x}$ branch while
$(N+1)!/(N+1+j)! = \prod_{i=2}^{j+1}(N+i)^{-1} \le (N+2)^{-j}$ gives the geometric
branch when $x < N+2$. \eqref{eq:bnd:kasm} follows because every entry of $G$ is a signed
sum of at most eight of the thirty-six scaled pairs.
\end{proof}

\paragraph{Measured (script \code{chkKser.jl}, $300$-bit).} Both the pointwise moment
inequality and the resulting remainder bound were checked on real face pairs of the
lattice, with the moments from the closed forms of the companion document and the
``truth'' the tail of the same series summed to $m_{\max} = 60$. At
$\mathbf{n} = (2,0,0)$ of the $\lambda/32$ cube the monotonicity
$|I_{n-1}| \le D_{\max}^{n}I_{-1}$ holds with a worst ratio of $0.0802$, and
$\mathrm{bound}/\mathrm{truth}$ is $4.44$, $8.05$, $14.4$, $25.2$, $43.2$ at
$N = 8,12,16,20,24$ for $f = 1$ and $4.46$, $8.07$, $14.4$, $25.3$, $43.3$ at
$f = 1+0.1\ii$; at $(3,1,0)$ of the same cell the ratios are $4.07$--$20.5$ over the
same $N$ and both frequencies, with monotonicity ratio $0.111$. On the slender cell,
which is the band the $k$-series actually owns, the offsets $(0,0,2)$, $(0,0,8)$ and
$(1,1,17)$ give $x = |k|D_{\max} = 0.280$, $0.299$, $0.593$, monotonicity ratios
$0.0198$, $0.0273$, $0.0595$, and $\mathrm{bound}/\mathrm{truth}$ between $7.48$ and
$266$ over $N = 8,12,16,20$ and both frequencies, with no violation in any of the $44$
rows. The bound therefore over-states the term count by at most three terms,
which matches the wider sweep on which it was originally measured
($N_{64}-N_{\mathrm{bnd}} \in [-3,+3]$ over $233$ rows). The assembly amplification
$\Lambda$ is the expensive part: $68.5$ at $(2,0,0)$ of the $\lambda/32$ cube, up to
$1405$ on its body diagonal, $15.4$ to $23.5$ at $\lambda/4$, and
$2.2\times10^{3}$ to $8.1\times10^{4}$ for the slender cell.

\subsection{Two radii, and which statement is about what}
\label{sec:bnd:radii}

The Cartesian organization of the same expansion --- family (a) of the registry, which
sums $\sum_{\boldsymbol{\alpha}}(\partial^{\boldsymbol{\alpha}+\mathbf{e}_a+\mathbf{e}_b}
g)(\mathbf{R})\mu_{\boldsymbol{\alpha}}/\boldsymbol{\alpha}!$ --- comes with a Cauchy
bound whose convergence criterion is $\sqrt{2}\,\rho < 1$, and it is important to be
precise about what that criterion is a statement about, because the two thresholds have
been conflated.

\begin{proposition}[the radius of the series]
\label{prop:bnd:rho}
Fix $\boldsymbol{\delta} \neq 0$ and let $\varphi(t) = F_{ab}(\mathbf{R}+t\boldsymbol{\delta})$
with $F_{ab} = (\partial_a\partial_b+\delta_{ab}k^2)g$ continued to complex argument
through $u = \sqrt{\mathbf{z}\cdot\mathbf{z}}$. Then $\varphi$ is analytic in the disc
$|t| < |\mathbf{R}|/|\boldsymbol{\delta}|$ and has singularities on its boundary.
Consequently the degree-graded Taylor series of $F_{ab}$ about $\mathbf{R}$ converges
absolutely at every $\boldsymbol{\delta}$ with $|\boldsymbol{\delta}| < |\mathbf{R}|$,
and the integrated series $\sum_q \int_D w\,c_q(\boldsymbol{\delta})$ converges
absolutely whenever $\rho = r_{\mathrm{d}}/|\mathbf{R}| < 1$.
\end{proposition}

\begin{proof}
$g$ is singular exactly where $\mathbf{z}\cdot\mathbf{z} = 0$. Along
$\mathbf{z} = \mathbf{R}+t\boldsymbol{\delta}$,
$\mathbf{z}\cdot\mathbf{z} = |\mathbf{R}|^2 + 2t(\mathbf{R}\cdot\boldsymbol{\delta})
+ t^2|\boldsymbol{\delta}|^2$ is a quadratic in $t$ whose discriminant
$4(\mathbf{R}\cdot\boldsymbol{\delta})^2 - 4|\mathbf{R}|^2|\boldsymbol{\delta}|^2$ is
non-positive by Cauchy--Schwarz, so its two roots are complex conjugates with product
$|\mathbf{R}|^2/|\boldsymbol{\delta}|^2$; hence
$|t_+| = |t_-| = |\mathbf{R}|/|\boldsymbol{\delta}|$ \emph{exactly}. The last statement
follows by taking $\beta$ with $\rho < \beta < 1$, applying the Cauchy estimate on
$|t| = \beta/\rho_{\boldsymbol{\delta}}$ with $\rho_{\boldsymbol{\delta}} =
|\boldsymbol{\delta}|/|\mathbf{R}| \le \rho$, and bounding
$\sum_q\int_D w|c_q| \le (\sup_D M)\,V_{\mathrm{t}}^{2}\sum_q(\rho/\beta)^{q}$, the
supremum being finite because the relevant circles stay at distance
$(1-\beta)|\mathbf{R}|$ from the singular set.
\end{proof}

\begin{remark}[$\rho < 1$ is about the series, $\sqrt2\,\rho < 1$ is about the bound]
\label{rem:bnd:two}
The Cauchy bound of family (a) estimates the Taylor coefficients on a \emph{complex ball}
of radius $r_{\mathrm{d}}$ centred at $\mathbf{R}$, because that is what a multivariate
Cauchy estimate needs. The distance from $\mathbf{R}$ to the null cone
$\{\mathbf{z}\cdot\mathbf{z} = 0\}$ in $\mathbb{C}^3$ is $|\mathbf{R}|/\sqrt2$, not
$|\mathbf{R}|$: with $\mathbf{w} = \mathbf{u}_r + \ii\mathbf{v}$ and
$\mathbf{a} = \mathbf{R}+\mathbf{u}_r$, $|\mathbf{z}\cdot\mathbf{z}|^2 =
(|\mathbf{a}|^2-|\mathbf{v}|^2)^2 + 4(\mathbf{a}\cdot\mathbf{v})^2$ vanishes first at
$|\mathbf{w}| = |\mathbf{R}|/\sqrt2$. So the ball-based bound exists only for
$\sqrt2\,\rho < 1$, and it is silent --- not false, silent --- in the range
$1/\sqrt2 \le \rho < 1$. Proposition~\ref{prop:bnd:rho} shows the series converges there
anyway, because the real difference box never leaves the real directions, along which the
singularity sits at $|t| = 1/\rho > 1$.

Measured (script \code{chkRay.jl}, $300$-bit): the two roots have
$||t_\pm| - |\mathbf{R}|/|\boldsymbol{\delta}||\,/\,(|\mathbf{R}|/|\boldsymbol{\delta}|)
\le 8.5\times10^{-91}$ at $(2,0,0)$, $(2,2,0)$ and $(3,0,0)$ of the $\lambda/32$ cube, at
$(2,0,0)$ of the $\lambda/4$ cube, and at $(0,0,16)$ of the slender cell. At the worst
real corner $\boldsymbol{\delta} = \mathbf{s}$ of the cubic $(2,0,0)$, where
$\rho = 0.866$ and $\sqrt2\rho = 1.2247 > 1$ so that the Cauchy criterion fails, the
Taylor coefficients on the circle $|t| = 1.097$ obey $|a_q| \le M\beta^{-q}$ with
$M = 3.53\times10^{5}$ ($|a_q| = 1418$, $1616$, $230$, $6.43$, $3.66\times10^{-3}$ at
$q = 8,16,32,64,128$), and the partial sums converge at $t = 1$ with relative errors
$0.595$, $8.97$, $1.879$, $5.05\times10^{-3}$, $1.75\times10^{-5}$ at the same $q$ ---
slowly, with an observed asymptotic ratio of $0.887$ approaching $\rho = 0.866$, but
convergently. At $(2,2,0)$, where $\sqrt2\rho = 0.866 < 1$, the same partial sums reach
$4.70\times10^{-13}$ at $q = 64$ and $6.91\times10^{-25}$ at $q = 128$. At the slender
$(0,0,16)$, $\rho = 1.416 > 1$, the radius is $0.706 < 1$ and the series genuinely
diverges at the corner.

\emph{Resolution of a disagreement between the two earlier analyses.} One of them records
$(2,0,0)$ as a divergent, asymptotic series that ``never becomes a convergent tail'', on
the evidence of a non-monotone \code{BigFloat} error sequence
($7.9\times10^{-12}$ at $p = 72$ falling to $2.2\times10^{-15}$ at $p = 120$); the other
records it as convergent with an effective ratio of about $0.73$, consistent with the
multipole form needing $L = 94$ for $6.6\times10^{-14}$. Proposition~\ref{prop:bnd:rho}
and the measurement above settle it: the series \emph{converges} at $(2,0,0)$, slowly and
non-monotonically, and the observed plateau is the interaction of a slow convergent tail
with \code{Float64}/\code{BigFloat} rounding, not divergence. What is genuinely
unavailable there is a Cauchy \emph{bound}, which is the practical reason that offset is
handed to the octant split in Section~\ref{sec:join} rather than to a longer whole-box
sum.
\end{remark}

\subsection{What is not proven}
\label{sec:bnd:notproven}

Everything above bounds a truncation. Nothing above bounds a rounding error, and no
bound in this document is a proof that the \code{Float64} evaluation of the truncated sum
attains the truncated accuracy. That gap is real and it is where the remaining risk of the
method sits. What exists instead is measurement, and it should be read as measurement.

\begin{itemize}[nosep,leftmargin=1.4em]
\item \textbf{The geometry table is not evaluable in \code{Float64} at all.} The
contraction of the integer solid-harmonic coefficients \eqref{eq:exp:rec} against the box
moments has $\sum|\text{term}|/|\text{result}| = 2.1\times10^{5}$ at $L = 20$,
$5.8\times10^{8}$ at $L = 40$ and $7.31\times10^{11}$ at $L = 56$ for the cubes, and up to
$1.90\times10^{12}$ over the twelve audit shapes --- almost $12$ decimal digits. This is a
\emph{measured} amplification, not a bounded one. The response is to build the table in
exact integer and rational arithmetic and to contract at $256$ bits, which removes the
question rather than answering it; the one-off cost is $14$--$18$~s and $0.7$~GB of
peak resident memory per cell shape (Section~\ref{sec:use:cost}).
\item \textbf{The per-offset sums are measured to be benign, not proved to be.} The
$l$-sum has $\sum_l|t_l|/|T_{ab}| \le 8.4$ over the whole $\lambda/32$ table and
$\le 67$ at $\lambda/4$ on the axial $(1,1)$ entry; the growth is linear in
$|k\mathbf{R}|$, $c_l \approx 0.67|k\mathbf{R}|$, so it reaches $1.8$ digits at
$|k\mathbf{R}| = 100$. The $n$-sum has $3.07$--$4.14$. The sub-box sum has $1.0$--$1.55$.
All four are measurements over the offsets tested.
\item \textbf{The special functions.} The upward recurrence for $h_l$ is measured
accurate to $3\times10^{-15}$ relative for $l \le 40$,
$|z| \in [0.3,300]$, $\operatorname{Im}z \in [0,30]$ --- including
$z = 300+30\ii$ where $j_0/h_0 = 6\times10^{25}$ and the recurrence still returns $h_{40}$
to $7.8\times10^{-17}$, because the contamination injected at $l = 0$ propagates as
$j_l/j_0$, which decreases, while $|h_l|$ increases. That mechanism is understood but not
proved here. Miller's downward recurrence returns $j_l$ to $6.8\times10^{-15}$ at worst.
The real spherical harmonics from the normalised Legendre recursion have absolute error
never above $6.2\times10^{-15}$ over $l \le 40$ and five directions.
\item \textbf{The phase seed.} For lattice offsets $\rho^2 = \sum(n_is_i)^2$ is exactly
representable, so the only \code{Float64} error in $e^{2\pi\ii f\rho}$ is
$2\pi f\rho\cdot\mathrm{ulp}(\rho)$; measured, that costs $2.15$ digits at
$|k\mathbf{R}| = 300$, $2.99$ at $1000$ and $3.35$ at $3000$, and a double-double phase
(twelve flops) removes it, recovering $2.5$ of the $3.35$ digits at
$|k\mathbf{R}| = 3000$. At $\lambda/4$ on a $128^3$ block $|k\mathbf{R}|$ reaches about
$350$, so this matters and is cheap.
\item \textbf{The end-to-end statement is empirical.} Against the $220$-bit reference, over
the $427$ cached tensors of Section~\ref{sec:ver:ref} --- four shapes, offsets from two to
$128$ cells, five frequencies --- the shipped library is accurate to $8.60\times10^{-15}$ in
max-norm and $9.34\times10^{-15}$ per entry, losing at most $1.59$ decimal digits to
rounding (Section~\ref{sec:ver:dig}). Over the audit's $60$ random (shape, offset,
frequency) cases the rounding error is at most $16.4\,\varepsilon\max_{ab}|T_{ab}|$, with an
$l$-sum amplification of at most $3.14\times10^{3}$ (Section~\ref{sec:bnd:cert}). Those are
the numbers that justify the method; they are not theorems.
\end{itemize}
